\documentclass[a4paper]{amsart}

\usepackage{color}
\usepackage[colorlinks=true, citecolor=red]{hyperref}  %needs to come before amsrefs package to work with \cite{}*{}
\usepackage{amssymb,latexsym}
\usepackage[alphabetic]{amsrefs}
\usepackage[mathscr]{eucal}
\usepackage{pdfsync}
\usepackage{graphicx}
\usepackage{epstopdf}
\DeclareGraphicsRule{.tif}{png}{.png}{`convert #1 `dirname #1`/`basename #1 .tif`.png}
\usepackage[all]{xy}
\usepackage{titletoc}
%\usepackage{txfonts}   %  \fint  for integral average

%\usepackage{showkeys}  %options notref  notcite 

%\input{rgb.tex}

%\textcolor{DarkRed}{text in DarkRed}


\CompileMatrices

\theoremstyle{plain}
\newtheorem{theorem}{Theorem}[section]
\newtheorem{lemma}[theorem]{Lemma}
\newtheorem{proposition}[theorem]{Proposition}
\newtheorem{corollary}[theorem]{Corollary}

\theoremstyle{definition}
\newtheorem{definition}[theorem]{Definition}%[section]


\theoremstyle{remark}
\newtheorem{defn}[theorem]{Definition} 
\newtheorem{example}[theorem]{Example}
\newtheorem{remark}[theorem]{Remark}
\newtheorem{discussion}[theorem]{Discussion}

%\numberwithin{equation}{section}

%\newcommand{\abs}[1]{\lvert#1\rvert}

\DeclareMathOperator{\Lip}{Lip}
\DeclareMathOperator{\dist}{dist}
\DeclareMathOperator{\expected}{\mathbb{E}}
\DeclareMathOperator{\prb}{Prob}
\DeclareMathOperator{\spt}{spt}
\DeclareMathOperator{\lap}{\Delta}
\DeclareMathOperator{\Tr}{Tr}

\raggedbottom

%\addtolength{\textwidth}{2cm}
%\addtolength{\oddsidemargin}{-1cm}
%\addtolength{\evensidemargin}{-1cm} 
%\addtolength{\topmargin}{-1cm} 
%\addtolength{\textheight}{2cm}

 
%%% BEGIN DOCUMENT
\begin{document}

Here is my interpretation of Patzschke's transfer operator on functions (in his notations): The random weights are denoted by $p_i(\omega)=p_i(\omega(\emptyset))$. (Note that his trees start at the empty set.) Then we get

$$
\mathcal{L}f(\omega,\sigma) =\sum_{i=1}^{N(\omega)} p_i(\omega(\emptyset)) \int f(W_i(\bar{\omega},\omega)), i\sigma)\, \mathbb{P}(d\bar{\omega})
$$
where $W_i(\bar{\omega},\omega)$ coincides with the 'labelled tree' $\bar{\omega}$ whose level one subtree starting at $i$ is replaced by the labelled tree $\omega$.


So, I would replace in your formula in the conditional expectation $w_i(\tilde{\omega})$ by $w_i(\omega)$, similarly for $N$ and then one can take them and the sum out of the expectation.

(M 27/10: This was wrong, I had used the wrong pressure function. So, your nice interpretation is the correct one.)

For the V-variable case I had already suggested at Canberra the model
\[
\tilde{\Omega}_V:=\{(\omega,\sigma)=(\omega_1,\ldots,\omega_V,\sigma_1,\ldots\sigma_V): \omega\in\Omega, \sigma_v\in\tilde{\omega_v}, v=1,\dots,V\}.
\]
In your old notations. $\tilde{\omega_v}$ is the set of all infinite paths in $\omega_v$ You suggested the new notation $\Sigma(\omega_v)$.

(Probably you did not copy it or it was on the blackbord.) Perhaps your semiproduct sign for the space $\tilde{\Omega}_V$ might be a bit misleading?

The shift operator in my notes was
$$\tilde{T}(\omega,\sigma)=(T_{(\sigma_1)_1}\omega_1,\ldots,\omega_v, T\sigma_1,\ldots,T\sigma_V)}$$
This should be the same as in your notations.


For later purposes we will also need the upper indices with the meaning

$Y^\tau(\omega):=Y(T_\tau\omega)$, $\tau\in\Sigma_*$, for several random variables $Y$.

Therefore we should carefully decide where to put the enumeration of the V-tuples, above or below. It is not yet clear to me, but you can convince me.

(JH 26/10: I think I agree with all the above --- and as we were discussing in Canberra)

I have not yet read the very last part of your notes, since I came back from the university later than I expected. I am aware that you have no time for discussion in the next weeks, but I will write again soon.

(JH: 26/10: I think the transfer operator on functions in the V-variable setting is a little different from what we were discussing?  This is Section 2.   Please note that I am referring to the second version of my email, i.e. the 4 page version.) 

(M 27/10: We did not yet discuss the transfer operator for the V-tuples at Canberra explicitly. In my opinion your definition is exactly that what we need.
The corresponding version of Patzschke's Lemma 4.3.2 remains valid. In his proof on page 33 for our special choice of the pressure function $\varphi$ integration w.r.t.\ $\mathbb{P}(d\bar{ \omega}) $ in your case may be replaced by that w.r.t. $dP_V(a)$. One could also work with Patzschke's notations interpreting the analogue of his $\mathbb{P}$ as the probability distribution induced by (3.3) and the following lines in your joint paper on dimension results.
But your version in the notes is easier to understand. Moreover, the random Gibbs measure on the cylinder sets of length $k$ for $V$-tuples should have a similar form as in the recursive construction: the product of the corresponding weights and the associated martingale limit shifted according to the cylinder sets. (Choose $h(\omega,\sigma):={\bf 1}_{\{\tau\}}(\sigma |k)$ for an admissible $V$-tuple $\tau$ of words of length $k$ and consider the limit in the dual operator version.)

(M27/10(2): We have to iterate $\lambda^{-1} \mathcal{L}^*$ on normed measures, where in our case
$\lambda=\int \sum_{\gamma^a\in\Gamma^a} w^a_\gamma\, dP_V(a)$.
`Iteration' of  $\lambda^{-1}\sum_{\gamma^a\in \Gamma^a} w^a_\gamma$
leads to the martingale $X_n$. In the random recursive case one uses $\lambda=1$.)

(M 27/10: This would also be a possibility to introduce the random Gibbs measure on $V$-tuples directly by means of their values on the cylinder sets, like in the recursive case. Then Perron-Frobenius would not be needed.) 





\end{document}
